

\begin{introduction}

\begin{section}{Background}

\end{section}

\begin{section}{Aims and Objectives}
\begin{subsection}{Research Questions}
The aim of PhD is to design and build tools that promote the use of data in civic participation and the appropriate tools and processes to support civic advocacy. For coherence and referencing in case study descriptions, I will reiterate the research questions below:
\begin{enumerate}
  \item What is the level of engagement with data and literacy around using data that local communities currently have?
  \item What does it mean to participate in civic action through data sharing and use?
    \begin{enumerate}
    \item What are the barriers of enabling data-driven participation?
    \item What are the the different ways of enabling data-driven participation?
    \item What are appropriate tools methods for accessing, interpreting and collaboratively making sense of data?
    \end{enumerate}
  \item What is the role of data in civic participation and advocacy?
    \begin{enumerate}
    \item How to enable effective use of data by communities?
    \item Does the process of citizen sensemaking and giving people voices through data shift the power relationship between the public and decision-makers, and strengthen the community?
    \end{enumerate}
\end{enumerate}
\end{subsection}

\begin{subsection}{Research Approach}
\end{subsection}

\end{section}

\begin{section}{Summary of Contributions}
\end{section}

\begin{section}{Thesis Outline}

The remainder of this thesis consists of seven chapters that reflect upon two case studies and an extended study of outcomes to investigate what role does data play in civic participation, advocacy and action. Analysis of the iterative design process, development and deployment of two accessible data systems has provided insights to conclude with a new framework for effective use of data in civic participation and advocacy that pinpoints steps for purposeful use of data and sustained engagement with it around civic concerns.
\newline

Chapter 1 explores the relevant literature for this work. Use of data within communities and in civic action explores a wide range of disciplines from technical fields like sensing networks, pervasive computing and data systems to human computer interaction and computer-supported collaborative work. Purpose of this chapter is to point out different uses of data in decision-making and the way citizens are involved or not involved in these data-driver processes.
\newline

Chapter 2 describes the overall approach and methodology for conducting the research and the lens which the case studies are analysed and reflected upon.
\newline

Chapter 3 presents the iterative design, development and deployment of SenseMyStreet, a sensor commissioning toolkit for communities that enables them to use scientific grade environmental monitors to collect data about their neighbourhood. Chapter describes the evolution of the toolkit, the analysis of lessons learned from setting it up and the results from particular case study with group called SPACEforHeaton, where citizens are using the data that is being collected to evidence an issue and influence policies linked to their local areas. 
% should I include all the groups? It seems that groups who had an aim for the data will be successful.
\newline

Chapter 4 focuses on additional sources of data to promote civic advocacy and presents case study two, the design and iterative development process of Data:In Place, a online map based system for accessing and interpreting open data from governmental sources to use them in local decision making.
\newline

Chapter 5 reflects on two case studies and goes deeper into the barriers and issues around data use and appropriation for local benefit. Through engagement activities and evaluation of the case studies an understanding of relevant stakeholders and needed processes for effective use of data are identified.
\newline

Chapter 6 starts off with describing the role of data in the context of knowledge and value creation linked to complex decision making processes. Moving away from the assumption that data is the 'silver bullet', meaning if there is data information, knowledge and wisdom reveals itself, to identifying the key principles of knowledge creation from data that leads to action. Chapter concludes with a framework for effective use of data that helps understand the gaps and 
needed resources resources to promote data use within the civil society to actively participate in renewed democratic processes. 
\newline

Finally, Chapter 7 concludes and assesses the success of the PhD through the primary aims and goals. It includes routes to future research and challenges with the built tools and processes, giving suggestions for the sustainability of these developed solutions.

\end{section}

\end{introduction}

